\documentclass[11pt]{article}

\usepackage[T1]{fontenc}
\usepackage{amsmath,amssymb,amsthm}
\usepackage[margin=1in]{geometry}
\usepackage{microtype}

\newtheorem{theorem}{Theorem}
\newtheorem{lemma}[theorem]{Lemma}
\newtheorem{corollary}[theorem]{Corollary}
\theoremstyle{remark}
\newtheorem{remark}[theorem]{Remark}

\newcommand{\Q}{\mathbb Q}
\newcommand{\Z}{\mathbb Z}
\newcommand{\Ztwo}{\mathbb Z_2}
\newcommand{\Ot}{\mathcal O}
\newcommand{\Ut}{\mathcal U}
\newcommand{\vt}{v_2}
\newcommand{\vp}{v_p}

\title{Integrality of the five-step product recurrence\\
       (MathOverflow 323963)}
\author{}
\date{}

\begin{document}
\maketitle

\begin{abstract}
We prove that the rational sequence defined by
\[
 a_{n+5}a_n=\prod_{i=1}^{4}(a_{n+i}+1),
 \qquad a_0=a_1=a_2=a_3=a_4=1,
\]
consists entirely of integers.  Odd primes are controlled by a valuation-atom
induction.  The prime~2 is treated by an alternating first integral and a
change of variables which reveals a fourteen-step $2$-adic return.
\end{abstract}

\section{Statement and initial values}

All terms are positive rational numbers, so the recurrence is well-defined.
The first terms are
\[
1,1,1,1,1,16,136,9316,43398586,941718655098991,\ldots.
\]

\begin{theorem}\label{thm:main}
For every $n\geq 0$, one has $a_n\in\Z$.
\end{theorem}

Put
\[
 b_n=a_n+1.
\]
Since $a_n\in\Q_{>0}$, it is enough to prove
\[
 v_p(a_n)\geq 0
 \qquad\text{for every prime $p$ and every $n$.}                 \tag{1}
\]

\section{Odd primes}

Fix an odd prime $p$.  Define integer-valued valuation atoms $t_n$ by
\[
 t_0=t_1=t_2=t_3=t_4=0,
 \qquad
 t_{n+5}=\vp(b_{n+5})-t_n\quad(n\geq0).                         \tag{2}
\]
Thus
\[
 \vp(b_{n+5})=t_{n+5}+t_n.                                     \tag{3}
\]

\subsection{The valuation identity}

\begin{lemma}\label{lem:odd-identity}
For every $n\geq0$,
\[
 \vp(a_{n+5})=t_{n+1}+t_{n+2}+t_{n+3}+t_{n+4}.                 \tag{4}
\]
\end{lemma}

\begin{proof}
The first five cases follow from
\[
\begin{aligned}
 a_5&=16,\\
 a_6&=8b_5,\\
 a_7&=4b_5b_6,\\
 a_8&=2b_5b_6b_7,\\
 a_9&=b_5b_6b_7b_8.
\end{aligned}                                                    \tag{5}
\]
Here the powers of~2 have zero $p$-adic valuation.

Suppose (4) holds at $n$.  Taking valuations in the recurrence at $n+5$
and using (3) gives
\[
\begin{aligned}
 \vp(a_{n+10})
 &=\sum_{j=6}^{9}\vp(b_{n+j})-\vp(a_{n+5})\\
 &=\sum_{j=6}^{9}(t_{n+j}+t_{n+j-5})
   -\sum_{j=1}^{4}t_{n+j}\\
 &=\sum_{j=6}^{9}t_{n+j}.
\end{aligned}
\]
This is (4) with $n$ replaced by $n+5$.  The five initial cases therefore
cover all residue classes modulo~5.
\end{proof}

\subsection{Nonnegativity and confinement}

We prove simultaneously
\[
 t_m\geq0                                                        \tag{6}
\]
and the confinement assertion
\[
 t_m>0\quad\Longrightarrow\quad
 t_{m-4}=t_{m-3}=t_{m-2}=t_{m-1}=0.                             \tag{7}
\]

For the initial check, the relevant values are
\[
\begin{array}{c|l}
m&b_m\\ \hline
5&17\\
6&137\\
7&9317=7\cdot11^3\\
8&43398587=29\cdot163\cdot9181\\
9&941718655098992=2^4\cdot37\cdot1172683\cdot1356497.
\end{array}                                                       \tag{8}
\]
The five displayed integers are pairwise coprime (equivalently, their odd
parts are pairwise coprime).  Hence a prime occurring in one of them occurs
in none of the preceding four.  Since $t_m=\vp(b_m)$ for $5\leq m\leq9$,
this proves both (6) and (7) through $m=9$.

Assume inductively that they hold through $m=k+9$.  Lemma
\ref{lem:odd-identity} already implies
\[
 \vp(a_{k+10})=t_{k+6}+t_{k+7}+t_{k+8}+t_{k+9}\geq0,
\]
and hence $\vp(b_{k+10})=\vp(1+a_{k+10})\geq0$ by the
non-Archimedean inequality.

Set $T=t_{k+5}$.  If $T=0$, then (2) immediately gives
$t_{k+10}=\vp(b_{k+10})\geq0$.  Suppose therefore that $T>0$.
Confinement at $k+5$ gives
\[
 t_{k+1}=t_{k+2}=t_{k+3}=t_{k+4}=0,
\]
so (4) shows that $\vp(a_{k+5})=0$.

Expanding the product of four successive $b$'s gives the exact identity
\[
\begin{aligned}
 b_{k+10}a_{k+5}
={}&b_{k+5}+a_{k+6}+a_{k+7}b_{k+6}\\
 &+a_{k+8}b_{k+6}b_{k+7}
  +a_{k+9}b_{k+6}b_{k+7}b_{k+8}.                               \tag{9}
\end{aligned}
\]
Every summand on the right has valuation at least $T$.  Indeed,
$\vp(b_{k+5})=t_k+T\geq T$.  Formula (4) puts the atom $T$ into
the valuation of each of $a_{k+6},\ldots,a_{k+9}$; all additional
$b$-factors have nonnegative valuation because their two atoms in (3)
are nonnegative.  Thus
\[
 \vp(b_{k+10})\geq T
\]
because $a_{k+5}$ is a $p$-adic unit.  By (2), $t_{k+10}\geq0$.

Finally, if $t_{k+10}>0$, then (3) and $t_{k+5}\geq0$ give
$\vp(b_{k+10})>0$.  Therefore
$a_{k+10}=b_{k+10}-1$ is a $p$-adic unit.  Equation (4) now yields
\[
 0=\vp(a_{k+10})=t_{k+6}+t_{k+7}+t_{k+8}+t_{k+9}.
\]
All four terms are nonnegative, so they are all zero.  This is the next
instance of (7), and the simultaneous induction is complete.

Combining (4) and (6) proves
\[
 \vp(a_n)\geq0
 \qquad(n\geq0)
\]
for every odd prime $p$.

\section{The prime 2}

The five initial numbers $b_0,\ldots,b_4$ are all even, so the preceding
atom construction does not give nonnegative atoms at $p=2$.  We instead
first reduce the recurrence by an exact alternating invariant.

\subsection{The alternating invariant}

Define
\[
 R_n=
 \frac{a_na_{n+2}a_{n+4}}
 {a_{n+1}a_{n+3}(a_{n+1}+1)(a_{n+3}+1)}.                        \tag{10}
\]

\begin{lemma}\label{lem:R}
One has
\[
 R_n=
 \begin{cases}
 \dfrac14,&n\text{ even},\\[2mm]
 4,&n\text{ odd}.
 \end{cases}                                                     \tag{11}
\]
\end{lemma}

\begin{proof}
Cancellation gives
\[
 R_nR_{n+1}
 =\frac{a_na_{n+5}}
 {(a_{n+1}+1)(a_{n+2}+1)(a_{n+3}+1)(a_{n+4}+1)}=1.
\]
Moreover $R_0=1/4$, because $a_0=\cdots=a_4=1$.  Hence the values
alternate as in (11).
\end{proof}

Write
\[
 \mathcal T(z)=\frac{z(z+1)}2,
 \qquad
 K_n=\begin{cases}1,&n\text{ even},\\16,&n\text{ odd}.
       \end{cases}                                                \tag{12}
\]
Since
\[
 a_{n+1}a_{n+3}(a_{n+1}+1)(a_{n+3}+1)
 =4\mathcal T(a_{n+1})\mathcal T(a_{n+3}),
\]
Lemma \ref{lem:R} is equivalent to
\[
 a_na_{n+2}a_{n+4}
 =K_n\mathcal T(a_{n+1})\mathcal T(a_{n+3}).                    \tag{13}
\]

\subsection{The transformed recurrence}

Define, for $n\geq0$,
\[
 x_n=\frac{a_na_{n+2}}{\mathcal T(a_{n+1})}.                    \tag{14}
\]
All quantities are positive and nonzero.  Multiplying (14) at $n$ and
$n+2$, and then using (13), gives
\[
 x_nx_{n+2}=K_na_{n+2}.                                         \tag{15}
\]
In particular,
\[
 a_{n+2}=\frac{x_nx_{n+2}}{K_n}.                                \tag{16}
\]

Eliminating the $a$'s from (14)--(15) gives an autonomous recurrence with
alternating coefficients:
\[
 x_{n+3}=
 \begin{cases}
 \displaystyle
 \frac{128x_nx_{n+2}(1+x_nx_{n+2})}
      {x_{n-1}x_{n+1}},&n\text{ even},\\[4mm]
 \displaystyle
 \frac{x_nx_{n+2}(16+x_nx_{n+2})}
      {512x_{n-1}x_{n+1}},&n\text{ odd}.
 \end{cases}                                                     \tag{17}
\]
For completeness, (17) follows by substituting
$a_j=x_{j-2}x_j/K_j$ into the definition of $x_{n+1}$ and cancelling
$x_{n+1}$.  The initial values are
\[
 x_0=x_1=x_2=1,
 \qquad x_3=16.                                                   \tag{18}
\]

\subsection{The fourteen-step return}

Let
\[
 \Ot=\Ztwo,
 \qquad
 \Ut=\Ztwo^\times.
\]
We use $2^r\Ot$ and $2^r\Ut$ also for negative integers $r$.

\begin{lemma}[Fourteen-step $2$-adic return]\label{lem:return14}
Let $x_0,x_1,x_2,x_3\in\Q_2^\times$ satisfy
\[
 x_0,x_1,x_2\in\Ut,
 \qquad x_3\in16\Ut,
 \qquad x_0\equiv x_2\pmod4.                                  \tag{19}
\]
Generate $x_4,\ldots,x_{17}$ by (17).  Then
\[
\begin{array}{c|cccccccccccccc}
r&0&1&2&3&4&5&6&7&8&9&10&11&12&13\\ \hline
x_r\in
&\Ut&\Ut&\Ut&2^4\Ut
&\Ot&2^4\Ot&2^2\Ot&2^2\Ut&2^{-2}\Ut&2^2\Ut
&2^2\Ot&2^4\Ot&\Ot&2^4\Ut.
\end{array}                                                       \tag{20}
\]
At the end of the block,
\[
 x_{14},x_{15},x_{16}\in\Ut,
 \qquad x_{17}\in16\Ut,
 \qquad x_{14}\equiv x_{16}\equiv-x_0\pmod4.                  \tag{21}
\]
In particular, the hypotheses (19) return after fourteen steps.
\end{lemma}

\begin{proof}
We give the complete finite calculation.  Write
\[
 x_0=u_0,\qquad x_1=u_1,\qquad x_2=u_2,\qquad x_3=16u_3,          \tag{22}
\]
where all $u_i$ are units and $u_2-u_0\in4\Ot$.  Normalize the
remaining terms according to (20)--(21):
\[
\begin{array}{c|cccccccccccccc}
r&4&5&6&7&8&9&10&11&12&13&14&15&16&17\\ \hline
x_r
&z_4&16z_5&4z_6&4z_7&z_8/4&4z_9&4z_{10}&16z_{11}
&z_{12}&16z_{13}&z_{14}&z_{15}&z_{16}&16z_{17}.
\end{array}                                                       \tag{23}
\]
Substitution in (17), with cancellation at each preceding line, gives the
following straight-line identities:
\[
\begin{aligned}
z_4&=\frac{u_1u_3(1+u_1u_3)}{2u_0u_2},\\
z_5&=\frac{u_2z_4(1+u_2z_4)}{2u_1u_3},\\
z_6&=\frac{(1+u_2z_4)(1+16u_3z_5)}{u_1},\\
z_7&=\frac{(1+16u_3z_5)(1+4z_4z_6)}{u_2},\\
z_8&=\frac{(1+4z_4z_6)(1+4z_5z_7)}{u_3},\\
z_9&=\frac{z_6z_8(1+z_6z_8)}{2z_5z_7},\\
z_{10}&=\frac{z_7z_9(1+z_7z_9)}{8z_6z_8},\\
z_{11}&=\frac{z_8z_{10}(1+z_8z_{10})}{2z_7z_9},\\
z_{12}&=\frac{2z_9z_{11}(1+4z_9z_{11})}{z_8z_{10}},\\
z_{13}&=\frac{z_{10}z_{12}(1+4z_{10}z_{12})}{2z_9z_{11}},\\
z_{14}&=\frac{2z_{11}z_{13}(1+16z_{11}z_{13})}{z_{10}z_{12}},\\
z_{15}&=\frac{z_{12}z_{14}(1+z_{12}z_{14})}{2z_{11}z_{13}},\\
z_{16}&=\frac{z_{13}z_{15}(1+z_{13}z_{15})}{2z_{12}z_{14}},\\
z_{17}&=\frac{z_{14}z_{16}(1+z_{14}z_{16})}{2z_{13}z_{15}}.
                                                                    \tag{24}
\end{aligned}
\]
Although some denominators in the second half of (24) need not be units,
all cancellations below are exact.  We now give the finite certificate in
the form used in the formal proof.

Put
\[
 \alpha=\vt(z_4),\qquad \gamma=\vt(1+u_2z_4),
\]
and define the three exceptional quotients
\[
 q_{10}=\frac{1+z_7z_9}{8z_6},\qquad
 q_{15}=\frac{1+z_{12}z_{14}}{z_{10}},\qquad
 q_{16}=\frac{1+z_{13}z_{15}}{2z_{12}}.                         \tag{25}
\]
We shall prove
\[
\begin{aligned}
 &z_4,z_5,z_6\in\Ot,\qquad z_7,z_8,z_9\in\Ut,\\
 &q_{10}\in\Ot,\qquad q_{15},q_{16}\in\Ut,\\
 &z_{13},z_{14},z_{15},z_{16},z_{17}\in\Ut,\\
 &z_{14}\equiv z_{16}\equiv-u_0\pmod4.                          \tag{26}
\end{aligned}
\]

First, a unit squared is congruent to \(1\) modulo \(8\), and the sum of
two \(2\)-adic units has positive valuation.  The first two lines of (24)
therefore give
\[
 \vt(z_5)=\alpha+\gamma-1,\qquad \vt(z_6)=\gamma.
\]
Here \(\alpha\geq0\).  If \(\alpha=0\), then \(\gamma\geq1\); if
\(\alpha>0\), then \(\gamma=0\) and \(\alpha\geq1\).  Hence
\(z_4,z_5,z_6\in\Ot\).

Set
\[
\begin{aligned}
A&=1+u_2z_4,& B&=1+16u_3z_5,\\
C&=1+4z_4z_6,& D&=1+4z_5z_7.
\end{aligned}
\]
The element \(B\) is a unit.  Since \(z_4,z_6\in\Ot\), the element
\(C=1+4z_4z_6\) is a unit as well.  The relevant identities are
\[
 z_6=\frac{AB}{u_1},\qquad
 z_7=\frac{BC}{u_2},\qquad
 z_8=\frac{CD}{u_3}.
\]
It follows successively that \(z_7\) is a unit, then that
\(D=1+4z_5z_7\) is a unit, and finally that \(z_8\) is a unit.  Moreover
\[
 E:=\frac{u_1u_3+ABCD}{z_4}\in\Ut,                               \tag{27}
\]
and the sixth line of (24) reduces exactly to
\[
 z_9=\frac{DE}{u_1u_3}\in\Ut.                                  \tag{28}
\]
To verify (27), if \(z_4\) is a unit then \(A\) is even, so the numerator
in (27) is a unit; division by the unit \(z_4\) preserves this.  If \(t=\vt(z_4)>0\), then
\[
 BCD\equiv1\pmod{2^{t+1}},\qquad A=1+u_2z_4.
\]
Writing \(p_0=u_1u_3\), the first line of (24) gives
\[
 p_0+1=\frac{2u_0u_2z_4}{p_0}.
\]
Consequently
\[
 p_0+ABCD
 \equiv
 u_2z_4\left(1+\frac{2u_0}{p_0}\right)
 \pmod{2^{t+1}},
\]
whose coefficient of \(z_4\) is a unit.  This proves (27)--(28).

We next prove that \(q_{10}\) is integral.  Put
\[
 s=\frac{2z_4A}{p_0},\qquad Q=BCD,\qquad L=\frac{Q-1}{s}.
\]
The preceding straight-line identities give
\[
 B=1+4u_2u_3s,\qquad
 C=1+2u_3sB,\qquad
 D=1+sBC,
\]
and hence, exactly,
\[
 L=4u_2u_3+2u_3B^2+(BC)^2.
\]
Thus \(L\equiv-1\pmod4\).  Also
\[
 p_0\equiv2z_4-1\pmod4,\qquad s\equiv2z_4A\pmod4.
\]
A direct factorization of the numerator of \(1+z_7z_9\) now gives
\[
 q_{10}=\frac{F\,u_1}{4u_2p_0^2B},
 \qquad
 F=u_0u_2+(p_0+2A)L+AsL^2.
\]
All factors outside \(F/4\) are units.  Using \(u_0\equiv u_2\pmod4\)
and the preceding congruences,
\[
 F\equiv2z_4(A^2-u_2-1)\pmod4.
\]
Since
\[
 A^2-u_2-1=u_2(2z_4+u_2z_4^2-1),
\]
the right side is divisible by \(4\): if \(\alpha>0\), the factor
\(2z_4\) already supplies two powers of \(2\); if \(\alpha=0\), then
both terms inside the final parenthesis are even.  Hence
\[
 \vt(F)\geq2,\qquad q_{10}\in\Ot.
\]

The same identities also give the returned residue of \(z_9\):
\[
 z_9=\frac{D}{p_0^2}
 \left(2u_0u_2+p_0u_2+2A^2L\right).
\]
Reduction modulo \(4\), using \(L\equiv-1\), gives
\[
 z_9\equiv-u_2\equiv-u_0\pmod4.
\]
The only apparent correction is \(2z_4(u_2-z_4)\); it is divisible by
\(4\), since either \(z_4\) is even or both \(u_2\) and \(z_4\) are odd.

For the second half of the block, write
\[
 r=q_{10},\qquad P=z_7z_9,
\]
and introduce
\[
\begin{aligned}
H&=1+Pr,\\
I&=1+2z_9rH,\\
J&=1+4\frac{z_9}{z_8}rHI,\\
K&=1+8\frac1{z_8}rHIJ.
\end{aligned}
\]
Since \(r\in\Ot\), the element \(H\) is integral and \(I,J,K\) are
units.  Exact simplification of (24) yields
\[
\begin{aligned}
z_{10}&=\frac{Pr}{z_8},&
z_{11}&=\frac{rH}{2},&
z_{12}&=\frac{HI}{z_7},\\
z_{13}&=\frac{IJ}{z_8},&
z_{14}&=\frac{JK}{z_9}.
\end{aligned}
\]
Define
\[
 L_{15}=\frac{HIJK-1}{r}.
\]
Successive expansion of \(I,J,K\) gives
\[
 L_{15}=P+2z_9H^2
 +4\frac{z_9}{z_8}(HI)^2
 +8\frac1{z_8}(HIJ)^2.
\]
The first summand is a unit and all the others have positive valuation,
so \(L_{15}\) is a unit.  The exact formula
\[
 q_{15}=\frac{z_8}{P^2}\,(8z_6+L_{15})
       =\frac{z_8}{P^2}\,\frac{P+HY}{r}.
\]
therefore proves \(q_{15}\in\Ut\).  The next normalized term satisfies
\[
 z_{15}=Kq_{15}\in\Ut.
\]
Moreover \(J\equiv K\equiv1\pmod4\), so
\[
 z_{14}=\frac{JK}{z_9}\equiv-u_0\pmod4.
\]

It remains to prove that \(q_{16}\) is a unit.  Put
\[
 Y=IJK,\qquad
 T=\frac{Y-1}{H},\qquad
 W=P+Y^2+PT.
\]
Exact expansion gives
\[
 T=2z_9r
 +4\frac{z_9}{z_8}rI^2
 +8\frac1{z_8}r(IJ)^2.
\]
The first summand has valuation \(\vt(r)+1\), while the other two have
valuation at least \(\vt(r)+2\) and \(\vt(r)+3\).  Thus
\[
 \vt(T)=\vt(r)+1.
\]
Since \(Y-1=HT\), \(Y\) is a unit, and \(Y+1\) has positive valuation,
\[
 \vt(Y^2-1)\geq\vt(r)+2.
\]
On the other hand, \(P+1=8z_6r\), so
\[
 \vt(P+1)\geq\vt(r)+3.
\]
Consequently \(\vt(P+Y^2)\geq\vt(r)+2\), whereas
\(\vt(PT)=\vt(r)+1\).  Therefore
\[
 \vt(W)=\vt(r)+1.
\]
Finally, exact cancellation in (24) gives
\[
 q_{16}=\frac{Wz_7}{2IP^2r}.
\]
The numerator and denominator have the same valuation
\(\vt(r)+1\), proving \(q_{16}\in\Ut\).

We still need the returned residue of \(z_{16}\).  Put
\[
 S=q_{15},\qquad M=z_{13}z_{15}=\frac{YS}{z_8}.
\]
Then
\[
 z_{16}=\frac{Mq_{16}}{z_{14}},
 \qquad q_{16}=\frac{1+M}{2z_{12}}.
\]
Because \(z_{14}\) is a unit and \(z_{14}^2\equiv1\pmod8\), it is enough
to prove
\[
 Mq_{16}\equiv1\pmod4.                                         \tag{29}
\]

Suppose first that \(\vt(r)>0\).  Then
\(P\equiv-1\pmod{16}\).  The exact formula for \(q_{15}\) gives
\[
 \frac{P+HY}{r}\equiv
 P+2z_9+4\frac{z_9}{z_8}\pmod8.
\]
This follows by substituting the displayed expansions of \(I,J,K\);
together with \(P\equiv-1\pmod{16}\), it gives
\(M\equiv1\pmod4\).  If
\(\vt(r)\geq2\), reduction of the exact formulas modulo \(4\) gives
\[
 \frac{M-1}{2}\equiv z_9+1\equiv z_7-1,\qquad
 z_{12}\equiv z_7.
\]
If \(\vt(r)=1\), the additional first-order contribution of \(I\) gives
\[
 \frac{M-1}{2}\equiv z_9+1+z_9rH,\qquad
 z_{12}\equiv z_7+z_9r.
\]
In both subcases,
\[
 1+\frac{M-1}{2}\equiv z_{12}\pmod4,
\]
and hence \(M(1+M)/2\equiv z_{12}\pmod4\).  This proves (29).

Now suppose that \(\vt(r)=0\).  Then \(r\) is a unit and \(H\) has
positive valuation.  The exact formula for \(q_{15}\) gives
\[
 M=\frac{Y(P+HY)}{P^2r}\equiv P\equiv-1\pmod4.
\]
Writing \(T=(Y-1)/H\) as above and setting
\[
 V=\frac12(P+Y^2+PT),
\]
the expansion of \(T\), together with \(P+1\equiv0\pmod8\), gives
\[
 V\equiv P(z_9r+2)\pmod4.
\]
Since
\[
 q_{16}=\frac{z_7V}{P^2rI},
\]
and \(P^2\equiv I\equiv1\pmod4\), this reduces to
\[
 q_{16}\equiv
 \frac{z_7(z_9r+2)P}{r}\equiv-1\pmod4.
\]  Thus
\(Mq_{16}\equiv1\pmod4\), proving (29) in the second case as well.
It follows that
\[
 z_{16}\equiv z_{14}\equiv-u_0\pmod4.
\]

For completeness, let
\[
 \delta=\vt(z_{10}),\qquad
 \epsilon=\vt(z_{11}),\qquad
 \zeta=\vt(z_{12}).
\]
Because \(P,z_8\) are units, the formula for \(z_{10}\) gives
\(\delta=\vt(r)\geq0\).  The next two formulas give
\[
 \epsilon=\delta+\vt(1+z_8z_{10})-1,\qquad
 \zeta=1+\epsilon-\delta.
\]
If \(\delta>0\), then \(\vt(1+z_8z_{10})=0\), so
\(\epsilon=\delta-1\geq0\) and \(\zeta=0\).  If \(\delta=0\), then
\(z_8z_{10}\) is a unit, so
\(\vt(1+z_8z_{10})\geq1\); hence \(\epsilon\geq0\) and
\(\zeta=1+\epsilon\geq1\).

The unit properties of \(q_{15}\) and \(q_{16}\) say precisely that
\[
 \vt(1+z_{12}z_{14})=\delta,\qquad
 \vt(1+z_{13}z_{15})=1+\zeta.
\]
Thus \(z_{13},z_{14},z_{15},z_{16}\) are units.  Finally,
\(z_{14}z_{16}\equiv1\pmod4\), so
\(\vt(1+z_{14}z_{16})=1\); the last line of (24) then makes \(z_{17}\)
a unit.  This establishes all of (26), hence (20)--(21).
\end{proof}

\subsection{Completion of the $2$-adic argument}

The initial transformed window (18) satisfies the hypotheses of Lemma
\ref{lem:return14}.  Because the length~14 is even, the parity in (17) is
unchanged when the lemma is iterated.  Hence, for all $n\geq0$,
\[
 \vt(x_n)\geq \ell_{n\bmod14},                                  \tag{30}
\]
where
\[
 (\ell_0,\ldots,\ell_{13})
 =(0,0,0,4,0,4,2,2,-2,2,2,4,0,4).                              \tag{31}
\]

By (16),
\[
 \vt(a_{n+2})
 =\vt(x_n)+\vt(x_{n+2})-\vt(K_n).                              \tag{32}
\]
For even $n$, $\vt(K_n)=0$; for odd $n$, $\vt(K_n)=4$.  Substitution of
(31) into the right side of (32) gives, as $n$ runs through the residue
classes modulo~14,
\[
 (0,0,0,4,2,2,0,0,0,2,2,4,0,0).                              \tag{33}
\]
Every entry is nonnegative.  Therefore
\[
 \vt(a_n)\geq0
 \qquad(n\geq0).                                                \tag{34}
\]

\section{Conclusion}

For every $n$, the positive rational number $a_n$ has nonnegative valuation
at every odd prime by Section~2 and at the prime~2 by (34).  Its reduced
denominator therefore has no prime divisor, so it is~1.  This proves
Theorem~\ref{thm:main}.  Every identity and valuation argument used above,
including the finite cancellation certificate (26), has also been checked
in Lean~4 with Mathlib.

\begin{remark}
The single negative entry $-2$ in the transformed valuation word (31) is
harmless.  Formula (16) recovers $a_{n+2}$ from a product of two transformed
terms, and the paired lower bounds in (33) are all nonnegative.  The
exceptional high spikes seen in the sequence $\vt(a_n)$ merely increase some
of the non-unit variables $z_4,z_5,z_6,z_{10},z_{11},z_{12}$; they do not
disturb the fourteen-step return tube.
\end{remark}

\end{document}






